%
%============= CHAPTER BREAK =================%
%
\chapter{Areteaus of Cappadocia (ca 140--160 CE)}%
\justifying%
\vspace*{3mm}%
%
\noindent{}A famous Greek medical author, Areteaus of Cappadocia wrote roughly around the mid-2nd century.\endnote{%
Tekiner 2015:508.} \\[1.5mm]

\noindent{}\textbf{βιαστής} / \textit{biastēs} / “violence” | “violent person” %
is a noun lemma which is conflated with the adjective βιαστός / %
\textit{biastos} / “violent” in TLG search results. %
The 335 TLG hits on a TLG search of either thus encompass both of these lemmas. %
The adjective lemma, as well as its antonym adjective ἀβίαστος / %
\textit{abiastos} / “non-violent”|“unforced”, are commonplace in Attic Greek.\endnote{%
The antyonym adjective ἀβίαστος also appears in Hellenistic Greek inscriptions, %
e.g., I. Aeg. Thrace E205, mid-2nd to early 1st BCE.} %
It can have a substantival sense, essentially an adjective functioning as a noun, %
i.e., “the violent one[s]”. %
The noun lemma βιαστής / \textit{biastēs} reflects a newer, %
modified spelling of the far earlier noun βῐᾱτάς / \textit{biatas} %
(18 TLG hits, starting with Alcmon and Pindar). %
Plutarch used the ancient form only once (\textit{De defectu oraculorum} 426C). %
Then it essentially disappeared, except in the \textit{Scholia on Pindar}, % 
where the clarifying comments confirm that the old noun had fallen out of usage.

The newer lemma first appears among non-Christian authors in Aretaeus, %
a passage we will analyze shortly. %
In early Christian texts, it would seem to appear first %
in Matthew 11:12.\endnote{%
Justin Martyrʼs \textit{Dial.Trypho} 51.3 %
only attests the adjective βιαστός. On the date of this text, %
see \textit{infra}, chapter on Valens, s.v., ἰσάγγελος.} %
This verse is especially interesting in that it contains a saying of Jesus with %
close parallels in Marcionʼs *Ev 16:16 and canonical Luke 16:16. %
A different, simpler version of this saying is quoted for Marcionʼs \textit{Evangelion} %
both by Tertullian in Latin and Epiphanius in Greek, and both lack this specific lemma.\endnote{%
“The law and the prophets until John and everyone forcibly enters it ”/ %
ὁ νόμος καὶ οἱ προφῆται ἕως Ἰωάννου καὶ πᾶς εἰς αὐτὴν βιάζεται %
(\textit{Pan}. 42.11.6 μγ (43); 42.11.17 Σχ. μγ (43); cf. 42.11.17
Ἔλ. μγ (43); GCS 31:113, 142); “Saying, ‘The law and the prophets were until John, %
from which the kingdom of god is announced’” / %
\textit{dicens lex et prophetae usque ad Ioannem ex quo regnum dei adnuntiatur} %
(\textit{Marc}. 4.33.7; SC 456:406, 408; Evans 446).} %
In subsequent Christian literature, the term is often repeated, %
but almost always in reference to the Matthean saying. %

In the passage from Areteaus, the term accentuates %
that, when afflicted with gout, women endure more intense pain than men. %

\vspace*{3mm}
 \begin{quote}γυναῖκεϲ δ’ ἀραιότερον μὲν ἀνδρῶν, χαλεπώτερον δέ· %
τὸ γὰρ μὴ ἐν ἔθεϊ μηδὲ οἰκεῖον, ἢν ὑπὸ ἀνάγκηϲ κρέϲϲον γένηται, %
\textbf{βιαϲτὴν} μέζονα τίκτει τὴν ξυμφορήν. (\textit{S.D.} 2.12.12)\endnote{%
TLG 0719.002, ed. Hude 1958.}% 
\\[1mm] %
 Though women have it more infrequently than men, it is worse. %
 For what is neither customary nor familiar, which may become necessarily greater, %
 births a most \textbf{violent} consequence. \end{quote} % 
\vspace*{3mm}

In this volume we do not assume that Justinʼs claim to quote a saying %
by Jesus establishes that saying as reflecting an earlier, stable text form of Matthew %
or even necessarily an oral tradition. %
Elsewhere in this volume and in our previous research (see esp. \textit{SGC}), %
we have surfaced evidence that Justin likely had a hand %
in the editing of canonical Matthew and canonical Mark. %
Here Justin quotes the saying for highly polemical purpose, to aver that Jesus Christ --- %
as \textit{the} decisive, new covenant ---, %
explicitly said that John the Baptist was the last Judean prophet, %
after whom no other Judean prophets would arise. %
In view of subsequent Christian history (from Montanists to Pentecostals), %
this strategy of charismatic suppression rarely proved successful. %

Justin thus has Jesus \textit{in voce propriae} relegate all %
legitimate prophecy to himself and his Judean/Israelite predecessors, %
while denying any legitimate future prophecy, even among his followers. %
The Marcionite and canonical Paulines, and canonical and non-canonical Acts, %
show this was quite a contentious issue for 2nd century Christian communities. %
Justinʼs dominical \textit{logion} here may reflect his own tampering with Matthew, %
and his patriarchal, scribal, early-Orthodox aversion toward %
Judean and/or Christian charismatic prophets (male and female). %

If the editor of canonical Matthew and/or Justin picked up this specific term %
from the medical author Areteaus %
(through whatever mode and across whatever degrees of separation and transmission), %
then such usage would be quite meaningful. %
The saying occurs within a confrontation between Jesus and Pharisees in Matthew, %
and within Justinʼs pioneering anti-Judean polemic. %
Both may imply Jerusalem was a disease-stricken woman whose pain (i.e., fall) %
the prophetic Jesus had foretold in dire, even medical terms. %
A borrowing from roughly contemporaneous medical terminology is interesting to consider, %
even if the dates of Aretaeus and clarity of dependence are uncertain. \\[1.5mm] %

\noindent{}\textbf{ἐλαφρία} / \textit{elaphria} / “levity|lightness|alleviation” appears just once in the NT, %
yet 105 times in total in the TLG. %
This noun is cognate to, yet distinct from, the adjective ἐλαφρός / \textit{elaphros} / %
“light in weight or movement, nimble” (2636 TLG hits, starting with Homer) %
and the verb ἐλαφρίζω / \textit{elaphridzō} / “to lighten” (131 TLG hits, starting with Archilochus, 7th c. BCE), %
among other variations. %
The first occurrence of the noun ἐλαφρία / \textit{elaphria} among extant and digitized non-Christian writings %
is found within Aretaeusʼ \textit{De curatione diuturnorum morborum} / %
\textit{On the Cure of Chronic Diseases}. %
In the section on the treatment of diabetes, %
as Aretaeus prepares to describe an involved treatment regimen, %
and after comparing diabetes and dropsy, %
he states (\textit{cur}. 2.2.1): %

\vspace*{2.5mm} 
 \begin{quote}ἥδε τοῖϲι ὑδεριῶϲι ἡ ὁδὸϲ ἐπιγίγνεται, %
 ἢν ἐϲ ἀγαθὸν ἡ νοῦϲοϲ τρέπηται· 
 ἀγαθὸν δὲ ἡ λύϲιϲ τῆϲ αἰτίηϲ καὶ μὴ τοῦ ἄχθεοϲ \textbf{ἐλαφρίη} μοῦνον. \\[1mm] %
 The path [to recovery] for those with dropsy has already been laid out, %
 by which the sickness is turned to good. %
 But what is good is the release of the cause and not merely the \textbf{alleviation} of the burden. %
 \end{quote} % 
 \clearpage

\noindent{}The term occurs once in the NT, %
in a part of 2Cor missing from Marcionʼs \textit{Apostolos}, %
one filled with gratuitous travel details and rhetorical fluff: %

\vspace*{3mm} 
 \begin{quote}Καὶ ταύτῃ τῇ πεποιθήσει ἐβουλόμην πρότερον πρὸς ὑμᾶς ἐλθεῖν, %
 ἵνα δευτέραν χάριν σχῆτε, %
 ¹⁶ καὶ δι’ ὑμῶν διελθεῖν εἰς Μακεδονίαν, καὶ πάλιν ἀπὸ Μακεδονίας
 ἐλθεῖν πρὸς ὑμᾶς καὶ ὑφ’ ὑμῶν προπεμφθῆναι εἰς τὴν Ἰουδαίαν. %
 ¹⁷ τοῦτο οὖν βουλόμενος μήτι ἄρα τῇ \textbf{ἐλαφρίᾳ} ἐχρησάμην; %
 ἢ ἃ βουλεύομαι κατὰ σάρκα βουλεύομαι, %
 ἵνα ᾖ παρ’ ἐμοὶ τὸ Ναὶ ναὶ καὶ τὸ Οὒ οὔ; (2Cor 1:15--17)\\[1mm] %
 And in this confidence I was planning previously to come unto you, %
 so that a second favor you may have, %
 ¹⁶ and through you to pass into Macedonia, %
 and again from Macedonia to come unto you, %
 and by you to be sent forth into Judea. %
 ¹⁷ Planning this, then, consequently did I conceive [this plan] \textbf{with levity}? %
 Or what I plan, do I plan based on flesh, %
 so that there be from me “yes, yes” and “no, no”?\end{quote} % 
 \vspace*{3mm}

By my lights, there is little else besides the new/rare lemma that %
connects canonical 2Cor to Athenaeus here, so it seems overmuch to postulate direct dependence. %
Even so, the use of the term in a mid-2nd century CE medical setting seems a %
far more plausible context of origin than %
a quasi-autobiographical segment of a letter by an Hellenistic Judean religious entrepreneur of the mid-1st century CE. %
The term also does not seem to appear in digitized Greek inscriptions, %
and is found only once in a Greek papyrus dated to the 6th century CE.\endnote{%
I.Cilicie 36 — Andana Mus. — 6th c. CE — Syria 2.1921.207,1 — cf. SEG 37.1340.} %
After Athenaeus and canonical 2Cor, its next attestation in the TLG is around the mid-4th century CE %
in the \textit{Homilies} of Symeon the Mesopotamian (commonly known as Ps-Macarius). %
The preponderance of evidence points to a mid-2nd century or later date for this portion of 2Cor. %